\section{Decisiones tecnológicas}
\label{sec:decisiones}

Usaremos \textbf{MongoDB} con índices \campo{2dsphere} y compararemos
\textbf{Google Open Buildings v3} con \textbf{Microsoft Building Footprints}.
Todo el flujo corre en Python dentro de Docker para que sea reproducible.

\begin{table}[H]
\centering
\small
\caption{Decisiones tecnológicas del proyecto.}
\label{tab:decisiones}
\begin{tabularx}{\textwidth}{@{}p{2.6cm}p{3.6cm}L@{}}
\toprule
\textbf{Decisión} & \textbf{Elección} & \textbf{Justificación} \\
\midrule
Motor NoSQL & MongoDB 8 (modelo documental) &
  Guarda geometrías en GeoJSON de forma nativa, tiene índices \campo{2dsphere}
  y operadores \campo{\$geoWithin} y \campo{\$geoIntersects}~\cite{mongodb}.
  Su pipeline de agregación con \campo{\$group} y \campo{\$merge} calcula los
  totales por municipio dentro de la base. Puede escalar horizontalmente con
  sharding. \\
Fuente de edificios 1 & Google Open Buildings v3~\cite{google} &
  Cubre América Latina con imágenes de 50 cm (inferencia de mayo de 2023).
  Trae \campo{area\_in\_meters} y \campo{confidence} (de 0,65 a 1,0), lo que
  permite validar nuestras áreas y hacer análisis de sensibilidad. \\
Fuente de edificios 2 & Microsoft Building Footprints~\cite{microsoft} &
  Imágenes de 2014 a 2021 con licencia ODbL. Está en GeoParquet particionado
  por país (\campo{RegionName}) y por quadkey de nivel 9, así que se puede
  descargar solo Colombia. Solo trae la geometría, por lo que el área se
  calcula. \\
Límites municipales & MGN 2025 del DANE, nivel municipio~\cite{mgn} &
  Es la fuente oficial que exige el enunciado. Se filtra con los 170 códigos
  de \campo{MunicipiosPDET.xlsx}. \\
Sistemas de referencia & Almacenamiento en WGS84 (EPSG:4326); áreas en
  MAGNA-SIRGAS Origen Nacional (EPSG:9377) &
  GeoJSON y \campo{2dsphere} exigen longitud y latitud en WGS84. Las áreas se
  calculan en una proyección plana oficial de Colombia para que queden en
  metros cuadrados. \\
Herramientas & Python 3.12, GeoPandas, Shapely, PyMongo, DuckDB, Docker
  Compose, GitHub &
  GeoPandas valida y reproyecta. DuckDB lee los Parquet y filtra antes de
  cargar. PyMongo carga por lotes. Docker Compose levanta la misma base en
  cualquier máquina. \\
\bottomrule
\end{tabularx}
\end{table}

% =====================================================================
% 3. Plan de implementación
% =====================================================================
